\subsection{Two Stores, Deliberately Separated}

Two kinds of data exist in the system, and almost every design decision in this section
follows from refusing to mix them.

The knowledge base is immutable. It is a serialised artefact produced by the pipeline of
Chapter~\ref{ch:pipeline}, shipped with the application, loaded into memory at startup, and
never written to at runtime. It lives in no database at all. The consequence worth stating
is that the knowledge base cannot be corrupted by application traffic: there exists no code
path by which a user action modifies what the system believes the document says. A user can
create an account, hold a hundred conversations and share every one of them, and the answer
to ``who approves 2.126'' is bit-for-bit what it was before they arrived.

Application state is the opposite in every respect. Accounts, conversations, and sharing
grants are mutable by definition, they are written on ordinary user actions, and they must
survive a process restart. They therefore live in a relational database.

Keeping the two apart is what makes the correctness argument of Chapter~\ref{ch:agent}
survive contact with a multi-user application. If the extracted matrix lived in the same
database as the conversations, then every feature added to the platform would become a
potential vector for altering the facts, and each one would need to be argued against that
possibility. Because it does not, the argument is made once, structurally, and does not need
revisiting.

\subsection{Choosing the Persistence Layer}

The database is a managed PostgreSQL instance hosted on Neon, reached through SQLAlchemy as
an object-relational mapper. Two aspects of that choice were forced rather than preferred,
and both are worth recording because the reasoning generalises.

The first is why a hosted database at all, rather than a file alongside the application. The
application is deployed on a platform whose free tier provides no persistent disk. A
file-based database written next to the running service would therefore be destroyed on
every redeployment, and the failure would be silent: the service would restart cleanly, the
schema would be recreated empty, and every registered account would simply be gone with no
error anywhere to indicate it. That is the same class of defect as the memory failure
described in Section~\ref{sec:deployment}, an assumption about the persistence of local
state that the platform does not honour, and it was avoided in advance only because the
earlier failure had already taught the lesson.

The second aspect is the local-development fallback. Requiring a hosted database to run the
test suite would make the tests slow, network-dependent, and impossible to run offline, so
the connection string defaults to a local file-based SQLite database when no environment
variable is supplied. Production sets that variable and gets PostgreSQL. This yields a
genuine engineering problem, since the two engines do not offer the same column types, and
the way it is resolved is described next.

\subsection{Identifiers and Dialect Portability}

Every table uses a randomly generated UUID as its primary key rather than an
auto-incrementing integer. The reason is not aesthetic. Conversation identifiers appear in
URLs, and a sequential integer key would make those URLs enumerable: a user holding
conversation 41 could try 40 and 42 and learn, from the difference between a refusal and a
not-found, how many conversations exist and roughly when theirs was created. Authorisation
is still checked on every request, so enumeration would not grant access, but a random
128-bit identifier removes the information leak entirely rather than relying on the check to
contain it. Identifiers are generated in the application rather than by the database, which
also means an object's identity is known before it is persisted.

The portability problem is that PostgreSQL has a native \texttt{UUID} column type and SQLite
does not. Rather than abandoning the native type in production, or maintaining two model
definitions, the schema declares the column once with a dialect variant: the type resolves to
a native \texttt{UUID} when the engine is PostgreSQL and to a 36-character string when it is
SQLite. One model definition drives both, so the objects exercised by the 375 tests on SQLite
are the same objects that run against PostgreSQL in production, and a schema change cannot
be applied to one and forgotten in the other.

\subsection{The Accounts Table}

The \texttt{users} table is where the three sign-in paths converge, and its most important
property is a column that is allowed to be empty. Table~\ref{tab:usersschema} gives its
definition.

\begin{longtable}{|p{0.22\textwidth}|p{0.16\textwidth}|p{0.52\textwidth}|}
\caption{The \texttt{users} table.}\label{tab:usersschema}\\
\hline
\rowcolor{myblue}\textcolor{white}{\textbf{Column}} & \textcolor{white}{\textbf{Type}} & \textcolor{white}{\textbf{Purpose and constraints}} \\
\hline
\endfirsthead
\hline
\rowcolor{myblue}\textcolor{white}{\textbf{Column}} & \textcolor{white}{\textbf{Type}} & \textcolor{white}{\textbf{Purpose and constraints}} \\
\hline
\endhead
\texttt{id} & UUID & Primary key, generated by the application. \\ \hline
\texttt{email} & string & Unique and indexed, never null. Normalised to lower case on write so that address casing cannot create two accounts. \\ \hline
\texttt{name} & string & Display name. Nullable, since a provider may not release it. \\ \hline
\texttt{hashed\_password} & string & \textbf{Nullable.} The bcrypt digest for accounts that have a password. Null for accounts created through a federated provider. \\ \hline
\texttt{oauth\_provider} & string & Nullable. Which provider authenticated this account, where one did. \\ \hline
\texttt{oauth\_sub} & string & Nullable. The provider's own immutable subject identifier for this person. \\ \hline
\texttt{role} & enum & \texttt{analyst} or \texttt{administrator}. Never null. \\ \hline
\texttt{created\_at} & timestamp & Account creation, in UTC. \\ \hline
\multicolumn{3}{|p{0.92\textwidth}|}{\textbf{Constraint} \texttt{uq\_oauth\_identity}: the pair (\texttt{oauth\_provider}, \texttt{oauth\_sub}) is unique, so one provider identity cannot be attached to two local accounts.} \\ \hline
\end{longtable}

A single table serves all three paths, and a row is distinguished only by which columns carry
values. A password account populates \texttt{hashed\_password} and leaves the two federated
columns null. A federated account does the reverse. The schema permits both at once, which is
what allows a person who signed up with a password to later authenticate through Google and
arrive at the same account rather than a second one.

The nullable password column is a security property rather than a convenience. An account
created through Google has no password anywhere in this system, so there is no digest to
extract from a database disclosure, no credential to reuse against another service, and no
password to be phished. The login endpoint enforces the corresponding rule explicitly: it
refuses when \texttt{hashed\_password} is null, so a federated account cannot be attacked
through the password form at all. Figure~\ref{fig:dbusers} shows the table as deployed, with
the bcrypt digests visible and the federated columns null for the three accounts registered
at the time of writing.

\begin{figure}[htbp]
\centering
\includegraphics[width=0.98\textwidth]{figures/db-users.png}
\caption{The \texttt{users} table in the deployed instance. Each account carries a bcrypt
digest beginning \texttt{\$2b\$12\$}, and the federated identity columns are null for
accounts registered with a password.}
\label{fig:dbusers}
\end{figure}

\subsection{The Conversations Table}

A conversation is stored as one row carrying its entire message sequence in a single JSON
column, rather than as a parent row with one child row per message. This is a deliberate
denormalisation and it deserves the scrutiny it usually attracts.

The normalised design is the textbook answer and would be correct if messages were ever
addressed individually. They are not. Every read in this application fetches a whole
conversation to display it, and every write appends to the end of one. There is no query
anywhere in the system that asks for a single message, no pagination within a conversation,
and no aggregate computed across messages. A child table would therefore contribute a join
on every read and an ordering column to maintain, in exchange for a capability nothing uses.

The second reason is that the provenance metadata attached to each answer, the resolved
identifier, the resolution method, the confidence score, the model that phrased it and the
language it was answered in, is a nested structure whose shape changed repeatedly during
development. In a column-per-field design each of those changes is a schema migration. In a
JSON column the client and the agent agree on the shape and the database stores what it is
given. Given that this metadata is only ever read back and rendered, never filtered or
joined on, the flexibility was worth more than the constraint.

The honest counterpart is stated plainly: this choice forecloses future queries such as
``how many answers fell back to the template'' without either a scan over the JSON documents
or a migration to a normalised form. That is an acceptable trade only because no such query
exists today, and it would need revisiting if the usage analytics discussed in
Section~\ref{sec:govkpi} were implemented.

Table~\ref{tab:convschema} gives the definition, and Figure~\ref{fig:dbconversations} shows
it populated.

\begin{longtable}{|p{0.22\textwidth}|p{0.16\textwidth}|p{0.52\textwidth}|}
\caption{The \texttt{conversations} table.}\label{tab:convschema}\\
\hline
\rowcolor{myblue}\textcolor{white}{\textbf{Column}} & \textcolor{white}{\textbf{Type}} & \textcolor{white}{\textbf{Purpose and constraints}} \\
\hline
\endfirsthead
\hline
\rowcolor{myblue}\textcolor{white}{\textbf{Column}} & \textcolor{white}{\textbf{Type}} & \textcolor{white}{\textbf{Purpose and constraints}} \\
\hline
\endhead
\texttt{id} & UUID & Primary key. Appears in share URLs, hence random rather than sequential. \\ \hline
\texttt{owner\_id} & UUID & Foreign key to \texttt{users}, indexed. The account that created the conversation and the only one that may extend it. \\ \hline
\texttt{title} & string & Display title, defaulting to ``New chat''. \\ \hline
\texttt{title\_is\_default} & boolean & Whether the title is still the placeholder. \\ \hline
\texttt{messages} & JSON & The ordered message list, each entry carrying role, text, and provenance metadata. \\ \hline
\texttt{created\_at} & timestamp & Creation time, UTC. \\ \hline
\texttt{updated\_at} & timestamp & Maintained automatically on every write; drives sidebar ordering. \\ \hline
\end{longtable}

The \texttt{title\_is\_default} flag solves a small problem neatly enough to be worth a
sentence. A conversation should take its title from the first question asked, so that the
sidebar is navigable, but a user who renames a conversation must not have that title
overwritten by the next message. Storing the title alone cannot distinguish a placeholder
from a deliberate choice that happens to resemble one. The boolean records intent rather
than inferring it: the first message replaces the title only while the flag is set, and any
explicit rename clears it permanently.

\subsection{The Sharing Tables}

Sharing is carried by two tables that do different jobs, and separating them is what keeps
authorisation coherent.

\texttt{conversation\_shares} records grants that already exist: one row per pair of
conversation and recipient. It is the single authority on who may view a conversation,
and every read path consults it. The pair is constrained unique, which makes re-sharing with
the same colleague idempotent rather than an error or a duplicate row. Two separate foreign
keys point at \texttt{users}, one for the recipient and one for the grantor, so the table
also records who extended the access rather than only that it exists.

\texttt{conversation\_share\_links} records an invitation that has not yet been accepted: a
random token, the conversation it refers to, who issued it, and when it expires, defaulting
to seven days. Crucially, a link is not itself a grant. Claiming one creates an ordinary
row in \texttt{conversation\_shares}, which means the QR-code path and the by-name path
converge on the same authorisation record, and the question of who may view a conversation
continues to have exactly one answer in the code. Had the link been treated as a credential
in its own right, every read path would have needed to check two conditions instead of one,
and the two would have drifted.

One active link is maintained per conversation. A second request for a QR code returns the
existing token rather than minting a new one, so a code already printed or already scanned
does not quietly stop working.

Table~\ref{tab:shareschema} gives both definitions, and Figure~\ref{fig:dbshares} shows a
live grant.

\begin{longtable}{|p{0.17\textwidth}|p{0.26\textwidth}|p{0.47\textwidth}|}
\caption{The two sharing tables.}\label{tab:shareschema}\\
\hline
\rowcolor{myblue}\textcolor{white}{\textbf{Table}} & \textcolor{white}{\textbf{Column}} & \textcolor{white}{\textbf{Purpose and constraints}} \\
\hline
\endfirsthead
\hline
\rowcolor{myblue}\textcolor{white}{\textbf{Table}} & \textcolor{white}{\textbf{Column}} & \textcolor{white}{\textbf{Purpose and constraints}} \\
\hline
\endhead
\texttt{conversation\_shares} & \texttt{id} & Primary key. \\ \hline
 & \texttt{conversation\_id} & Foreign key, indexed. \\ \hline
 & \texttt{shared\_with\_user\_id} & Foreign key, indexed. The account granted read access. \\ \hline
 & \texttt{shared\_by\_user\_id} & Foreign key. The account that granted it. \\ \hline
 & \texttt{created\_at} & When access was granted. \\ \hline
 & \multicolumn{2}{p{0.73\textwidth}|}{\textbf{Constraint} \texttt{uq\_conversation\_share}: (conversation, recipient) unique, making a repeat share idempotent.} \\ \hline
\texttt{conversation\_share\_links} & \texttt{id} & Primary key. \\ \hline
 & \texttt{conversation\_id} & Foreign key, indexed. \\ \hline
 & \texttt{token} & Unique and indexed. The random secret encoded into the QR image. \\ \hline
 & \texttt{created\_by\_user\_id} & Foreign key. Who issued the link. \\ \hline
 & \texttt{created\_at} & Issue time. \\ \hline
 & \texttt{expires\_at} & Expiry, defaulting to seven days after issue. \\ \hline
\end{longtable}

\begin{figure}[htbp]
\centering
\includegraphics[width=0.98\textwidth]{figures/db-shares.png}
\caption{The \texttt{conversation\_shares} table, showing a live grant. The row names the
conversation, the account granted access, and the account that granted it.}
\label{fig:dbshares}
\end{figure}

\subsection{Referential Integrity and Deletion}

Both sharing tables are declared as dependants of the conversation with delete-orphan
cascade, so removing a conversation removes its grants and its outstanding links in the same
operation. This matters more than ordinary tidiness. Without the cascade, deleting a
conversation would leave grant rows pointing at an identifier that no longer resolves, and an
outstanding share link would remain valid in the sense that its token still existed and had
not expired. Access-control records that outlive the thing they protect are precisely the
kind of residue that later produces a security incident, so the lifetime of a grant is bound
to the lifetime of its conversation at the schema level rather than in application code that
might be bypassed.

One implementation detail is worth recording because it caused a real error. The share table
holds two foreign keys into \texttt{users}, and the mapper cannot infer which one a given
relationship traverses when more than one candidate exists; both relationships therefore
declare their join column explicitly. The error surfaces at mapper configuration time rather
than at query time, which is the preferable failure, but it is opaque enough on first
encounter to be worth naming here.

\section{Access Control}

\subsection{Why Authentication Was Necessary}

The first version of this system had no accounts. Anyone reaching the service could ask
questions, and conversations lived in the browser's own storage. That was adequate while the
system did one thing, and it stopped being adequate the moment conversations needed to be
persisted and shared, for a reason that is structural rather than incidental: sharing
requires ownership, ownership requires identity, and identity requires authentication. There
is no way to grant a colleague access to something that belongs to nobody.

A second reason is specific to this subject matter. The analytics layer aggregates the
delegation structure of the institution, which is organisationally sensitive in a way that a
single lookup is not, and restricting it requires being able to tell who is asking.

\subsection{Password Storage}

Passwords are never stored in a form from which the original can be recovered. Each account
holds a bcrypt digest, produced through the passlib library, and verification re-derives a
digest from the submitted password and compares the two. Bcrypt was chosen over a
general-purpose cryptographic hash for reasons worth setting out, since the distinction is
frequently elided.

A hash such as SHA-256 is designed to be fast, which is the correct property for verifying
file integrity and precisely the wrong one for storing passwords: speed is exactly what an
attacker holding a stolen database wants, since it lets them test billions of candidate
passwords per second on commodity hardware. Bcrypt is deliberately slow and, more
importantly, adjustably slow. Its cost parameter sets the number of key-expansion rounds as a
power of two, and raising it by one doubles the work required both for a legitimate
verification and for every guess an attacker makes. The digests in this system carry a cost
of twelve, meaning $2^{12}$ rounds, which is visible in the \texttt{\$2b\$12\$} prefix of
every stored value in Figure~\ref{fig:dbusers}. The parameter is stored alongside the digest
rather than fixed in code, so the cost can be raised as hardware improves and existing
accounts continue to verify against their original setting until their next successful login.

Bcrypt also salts each password individually and stores the salt in the same field. The
consequence is that two accounts choosing the same password hold entirely different digests,
which defeats both precomputed rainbow tables and the simpler attack of noticing that two
rows share a value and therefore share a password.

The interaction between password storage and federated accounts is the property described
earlier: an account authenticated through a provider carries no digest at all, and the login
endpoint refuses it before any comparison is attempted.

\subsection{Token Design}

Authentication is stateless. A successful sign-in returns a signed JSON Web
Token\citep{rfc7519}, and the client presents it in the \texttt{Authorization} header of
every subsequent request. Table~\ref{tab:jwtclaims} lists what the token asserts.

\begin{longtable}{|p{0.14\textwidth}|p{0.30\textwidth}|p{0.46\textwidth}|}
\caption{Claims carried by the access token.}\label{tab:jwtclaims}\\
\hline
\rowcolor{myblue}\textcolor{white}{\textbf{Claim}} & \textcolor{white}{\textbf{Value}} & \textcolor{white}{\textbf{Why it is present}} \\
\hline
\endfirsthead
\hline
\rowcolor{myblue}\textcolor{white}{\textbf{Claim}} & \textcolor{white}{\textbf{Value}} & \textcolor{white}{\textbf{Why it is present}} \\
\hline
\endhead
\texttt{sub} & Account UUID & Identifies the caller. Every ownership check resolves from this. \\ \hline
\texttt{email} & Address & Lets the interface show the signed-in account without an extra request. \\ \hline
\texttt{role} & \texttt{analyst} or \texttt{administrator} & Carries the authorisation decision, avoiding a lookup on every request. \\ \hline
\texttt{iat} & Issue time & Establishes token age. \\ \hline
\texttt{exp} & Issue time plus eight hours & Bounds the damage a leaked token can do. \\ \hline
\end{longtable}

The token is signed HS256\citep{rfc7515}, a keyed signature using a secret held only by the
server. Signing establishes integrity, not confidentiality: the claims are encoded, not
encrypted, and anyone holding the token can read them. Nothing secret is therefore placed in
it. What signing guarantees is that a client cannot alter a claim, and the claim that matters
is \texttt{role}: without a signature, elevating oneself to administrator would be a matter
of editing a string.

The signing secret is read from the environment with no fallback default, so a deployment
that has not configured one fails immediately at startup rather than silently signing tokens
with a value an attacker could guess from the source. Failing loudly at boot is the correct
behaviour for a misconfiguration of this class.

Expiry is eight hours, chosen to span a working day. The trade-off is explicit: a shorter
window would force repeated sign-ins during ordinary use, and a longer one would leave a
usable credential in a browser long after the person had finished. Because verification is
stateless, a token cannot be revoked before it expires, which is the cost of not consulting
the database on every request. For this application that cost is acceptable; a system holding
higher-value operations would need a revocation list and the lookup that comes with it.

On the client the token is held in session storage rather than local storage. Both are
readable by any script executing on the page, so the distinction is not about resisting
script access; it is about lifetime. Session storage is cleared when the tab closes, whereas
local storage persists indefinitely, leaving a valid credential on a shared or unattended
machine long after the user believed they had finished. The cost is that closing the tab
requires signing in again, which was judged the right trade for a system holding governance
information.

\subsection{Federated Sign-In}

Two federated paths exist alongside the password form, through Google and through Microsoft.
The motivation is institutional rather than technical: staff at an organisation of this kind
already hold a corporate directory account, and requiring a separate password adds a
credential to be managed, forgotten, reused, or written down, while removing the ability of
the institution to revoke access centrally when someone leaves.

Both providers are integrated through OpenID Connect rather than by hand-coded calls to
provider endpoints. The client is configured with a discovery document URL, and the endpoints,
signing keys and supported parameters are retrieved from the provider at runtime. This matters
for durability: a provider rotating its signing keys or relocating an endpoint does not break
the integration, because nothing about those details is hardcoded. The requested scope is the
minimum the application needs, namely the identity token plus the address and display name.
The Microsoft integration additionally accepts a tenant identifier, defaulting to the
multi-tenant endpoint, so that the same build can be pointed at a single institutional
directory by configuration alone.

Both providers are registered only when their credentials are present in the environment. An
instance configured with neither still starts and still serves password sign-in, which keeps
local development and the test suite free of any dependency on external identity services.

\subsection{Resolving a Federated Identity to an Account}

The step where a provider's response becomes a local account is the most security-sensitive
part of the flow, and the order of the lookups is the substance of it.

The account is located first by the pair of provider and subject identifier, the provider's
own immutable handle for that person, and never by email address alone. This ordering is
deliberate. Email addresses are not stable identifiers: an organisation may reassign a
departed employee's address to a new hire, and a person may change their own address at the
provider. A system matching on email alone would hand the new holder of a recycled address
the previous holder's account and every conversation in it. Matching on the subject
identifier makes that impossible, because the identifier is never reissued.

Only when no account carries that provider identity does the resolution fall back to email,
and then only to link the provider identity onto an existing account rather than to
authenticate against it. This is what lets a user who registered with a password later sign
in through Google and arrive at the same account rather than a duplicate. The account is
created fresh only when both lookups fail.

The uniqueness constraint on the provider-and-subject pair enforces the invariant at the
database level, so a defect in this logic cannot produce two local accounts claiming the same
provider identity.

\subsection{Roles and Authorisation}

Two roles are defined. An \textbf{analyst} may ask questions, attach documents, and own and
share conversations. An \textbf{administrator} may additionally open the analytics dashboard.
The distinction reflects the sensitivity difference described earlier: an individual
conversation reveals what one person asked, whereas the dashboard aggregates the institution's
delegation structure.

Authorisation is enforced server-side as a dependency on the protected routes rather than by
hiding controls in the interface. An analyst who constructs the request by hand is refused
exactly as one who clicks, which is the only enforcement that means anything; a hidden button
is a usability affordance, not a control.

The first account created on an empty database becomes an administrator, which solves the
bootstrap problem of needing an administrator before one can be appointed. This is recorded as
a deliberate convenience rather than a complete design: promoting a second administrator still
requires a direct database edit, and a production deployment would want an administrative
interface for it.

\subsection{Hardening the Sign-In Path}

Three further protections apply specifically to authentication, and
Table~\ref{tab:threatmodel} sets out what each defends against.

Failed attempts are rate-limited by source address within a rolling five-minute window, which
raises the cost of credential stuffing from free to impractical without locking out a
legitimate user who has mistyped a password twice.

A failed sign-in returns one generic message regardless of whether the account exists. The
temptation to be helpful here, by saying that no account was found, produces an oracle that
converts a list of institutional email addresses into a list of confirmed users. The
implementation is careful on a related point: the same generic refusal is returned when an
account exists but has no password because it is federated, so the response does not
distinguish that case either.

A minimum password length of eight characters is enforced at registration. This is a weak
control on its own and is stated as such; it raises the floor without pretending to establish
password quality.

\begin{longtable}{|p{0.30\textwidth}|p{0.56\textwidth}|}
\caption{Authentication threats and the corresponding mitigation.}\label{tab:threatmodel}\\
\hline
\rowcolor{myblue}\textcolor{white}{\textbf{Threat}} & \textcolor{white}{\textbf{Mitigation in this system}} \\
\hline
\endfirsthead
\hline
\rowcolor{myblue}\textcolor{white}{\textbf{Threat}} & \textcolor{white}{\textbf{Mitigation in this system}} \\
\hline
\endhead
Database disclosure exposing passwords & Only salted bcrypt digests at cost twelve are stored; federated accounts hold no digest at all. \\ \hline
Offline brute-force against stolen digests & Bcrypt's cost parameter makes each guess expensive, and the parameter can be raised without a migration. \\ \hline
Rainbow-table attack & Per-password salts make precomputation useless and hide the fact that two accounts share a password. \\ \hline
Account enumeration & One generic refusal for unknown account, wrong password, and federated account alike. \\ \hline
Credential stuffing & Rate limiting per source address over a rolling window. \\ \hline
Privilege escalation by editing a token & The role claim is covered by the HS256 signature; any alteration invalidates the token. \\ \hline
Credential left on a shared machine & Tokens held in session storage, discarded when the tab closes, and expiring after eight hours regardless. \\ \hline
Conversation URL enumeration & Random UUID identifiers rather than sequential integers, with the ownership check retained as well. \\ \hline
Recycled email address inheriting an account & Federated identities resolved by immutable provider subject identifier before email is ever consulted. \\ \hline
Stale access after a conversation is deleted & Grants and outstanding links removed by database cascade with the conversation itself. \\ \hline
Misconfigured deployment signing with a default secret & The signing secret has no default; a deployment without one fails at startup. \\ \hline
\end{longtable}

